\documentclass{ximera}
\title{Tangent vectors}

\begin{document}

    \begin{abstract} Different definitions of tangent vectors.    \end{abstract}
    \maketitle

            A reference for this material is Chapter 3 of John M. Lee. Introduction to smooth manifolds. Second edition. Grad. Texts in Math., 218. Springer, New York, 2013. xvi+708pp. ISBN: 978-1-4419-9981-8.

    
    Let $M$ be a smooth $n$-dimensional manifold and $p \in M$. 

    \begin{definition}[Tangent space]
    The \emph{tamgenet space to $M$ at $p$} is defined as the set of linear functions $L : C^{\infty } (M) \to \mathbb{R} $ such that $L(fg) = L (f) g(p) + f(p)L(g)$ for all $ f,g \in C^{\infty}(M) $.         Elements of $T_pM$ are called the \emph{derivations at $p$} or \emph{tangent vectors at $p$}.
    \end{definition}

    \begin{proposition}[Curves define derivations]\label{prop:curve-to-derivation}
        Let $\gamma : I \to M$ be a smooth curve with $0 \in I$ and $\gamma (0) = p$. The map $L_{\gamma} : C^{\infty}(M) \to \mathbb{R}$ given by
        \[       L_{\gamma} (f) : = (f \circ \gamma )' (0)                  \]
        is a derivation at $p$.

    \begin{solution}
        Exercise.
        %$L_{\gamma}$ is linear because taking derivatives is linear. The Leibniz rule follows from the classical one. 
    \end{solution}
    \end{proposition}

    \begin{theorem}[Equivalence of definitions]\label{thm:tpm-definitions}
        The map $\gamma \mapsto L_{\gamma}$ defined in Proposition \ref{prop:curve-to-derivation} is a surjective function from the set of smooth curves passing through $p$ at $0$ to $T_pM$. Moreover, $T_pM$ is an $n$-dimensional vector space.
    \end{theorem}

    The proof will consist of several steps. Fix $L \in T_pM$. 

    \begin{proposition}[Derivative of constant]
        For a constant function $c \in C^{\infty}(M)$, one has $L (c) = 0$. 

    \begin{solution}
        \begin{align*}
            L(c) & = L ( 1 \cdot c) \\
            & = L (1) c + 1 \cdot  L (c) \\
            & =  L (c) +  L (c) \\
            & = 2 L (c) .
        \end{align*}
        This is only possible if $L(c) = 0$.
    \end{solution}
    \end{proposition}
    
    \begin{proposition}[Locality] \label{prop:locality-of-derivatives}
        If $f,g \in C^{\infty} (M)$ agree in a neighborhood of $p$, then $L(f) = L (g)$.
    \begin{solution}
        Let $U \subset M$ be an open neighborhood of $p$ where $f$ and $g$ agree. Let $\rho \in C^{\infty} (M)$ be a bump function around $p$ supported inside of $U$. Then 
        \[ f-g \equiv (f -g) ( 1 - \rho ) \] 
        (check this!), and
        \begin{align*}
            L( f) - L (g) & = L (f-g) \\
            & = L ((f - g) ( 1 - \rho ) ) \\
            & = L (f-g) ( 1 - \rho (p) )  + (f(p) - g(p)) L(1 - \rho) \\
            & = L(f-g) (0) - (0) L (1 - \rho ) \\
            & = 0. 
        \end{align*}
    \end{solution} 
    \end{proposition}

    \begin{proposition}[Restriction of derivations]\label{prop:restriction-of-derivatives}
        Let $ U \subset M$ be open containing $p$. For each $f\in C^{\infty}(U)$, define 
        \[   L^U (f) : = L (\tilde{f}) ,     \]
        where $\tilde{f} \in C^{\infty} (M) $ is a function such that $f = \tilde{f} $ in a neighborhood of $p$. This is a well-defined derivation $L^U  \in T_p U $.

    \begin{solution}
        For $f \in C^{\infty} (U)$, the function $\tilde{f} \in C^{\infty}(M)$ can be obtained by multiplying $f$ by a bump function around $p$ supported inside of $U$. Proposition \ref{prop:locality-of-derivatives} implies that $L^U (f) $ does not depend on the choice of $\tilde{f}$. The linearity and Leibniz rule for $L^U$ follow from the corresponding properties for $L$.  
    \end{solution} 
    \end{proposition}

    \begin{proposition}[Restriction of derivations is injective] The map $L \mapsto L^U$ is a linear isomorphism from $T_pM$ to $T_pU$.
    
    \begin{solution}
 %       If $L(f) \neq 0$, then $L^U (f\vert _U ) = L (f) \neq 0$. This means 
  %      \[ [ L \neq 0 ] \Rightarrow [L^U \neq 0]. \] 
        For any $D \in T_pU$, one can define $\tilde{D} \in T_pM$ by
        \[      \tilde{D} (f) : = D (f \vert _ U ).                 \]
        It is straightforward to show that $D \mapsto \tilde{D}$ is the inverse of $L \mapsto L^U$.
    \end{solution}
\end{proposition}


    \begin{remark}
        For a chart $(U , \varphi ) $ and $f \in C^{\infty}(U)$, by an abuse of notation, we often identify $f$ with $f \circ \varphi$. Because of this, we sometimes  write $\partial_i f$ instead of $\partial _i (f \circ \varphi ^{-1} ) \circ \varphi $. 
    \end{remark}


    \begin{proof}[Proof of Theorem \ref{thm:tpm-definitions}]  
        Let $(U, \varphi)$ be a chart with $\varphi (p) = 0$ and $\varphi (U) = B_1(0)$.  For each $i \in \{ 1, \ldots , n \}$, notice that $\partial _i \vert _p : C^{\infty}(U) \to \mathbb{R}$ given by 
        \[      \partial_i \vert _p (f) : =  \partial _i  f  (p)      \]
        is a derivation at $p$. Moreover, $\partial_i \vert _p = L _{\gamma_i}$, where $\gamma_i (t) = \varphi ^{-1} ( t e_i)$. We will show that 
        \begin{equation}\label{eq:tpm-basis}
              \{   \partial_1 \vert _p , \ldots,   \partial_n \vert _p    \}      
        \end{equation}
        is a basis of $T_pM$. 
        
        Fix $L \in T_pM$, and $L^U \in T_pU$ given by Proposition \ref{prop:restriction-of-derivatives}. We claim that 
        \begin{equation}\label{eq:derivation-is-partials}
            L ^U = \sum_{i = 1} ^ n  L ^U ( x_i )\,  \partial _i \vert _p . 
        \end{equation}
       To see that, fix $f \in C^{\infty } (U)$. By the Fundamental Theorem of Calculus, for any $x \in B_1(0)$ we have
       \begin{align*}
           f(x) & = f(0) + \int _0 ^1 \frac{d}{dt} (f (tx)) \, dt \\
           & = f(0) + \int_0 ^1 \nabla f (tx) \cdot x \, dt \\
           & = f(0) + \sum_{i = 1}^n x_i \int_0^1  \partial _i f(tx) \, dt \\
           & = f(0) + \sum_{i = 1}^n x_i  h_i (x)
       \end{align*}
       for some functions $h_i \in C^{\infty} (U)$ with $h_i (0) = \partial_if (0)$.  Finally, 
    \begin{align*}
        L ^U (f) & = L^U (f(0)) +\sum_{i = 1}^n L ^U \left(  x_i h_i    \right) \\
        & = 0 + \sum _i ^n \left[   L ^U( x_i ) \, h_i (0)  + 0   \cdot  L^U (h_i)    \right] \\
         &= \sum _{i = 1} ^n L^U (x_i) \, \partial_i  f (p).
    \end{align*}
    This proves \eqref{eq:derivation-is-partials}.     To see that \eqref{eq:tpm-basis} is linearly independent, assume  
    \[ \sum_{i = 1} ^n c_i \partial_i \vert _p  = 0\] 
    for some scalars $c_1, \ldots , c_n$. Then
    \begin{align*}
        0 & =  \sum_{i = 1} ^n c_i \partial_i \vert _p \left( \sum_{j =1 }^n c_j x_j \right)  \\
        & =  \sum_{i = 1} ^n  \sum_{j =1 }^n c_i c_j  \partial_i \vert _p x_j  \\
        & = \sum_{i = 1} ^n  \sum_{j =1 }^n c_i c_j \delta_{ij} \\
        & = \sum _{i = 1 }^n c_i ^2 .
    \end{align*}
    This would mean $c_1 = \ldots = c_n = 0$. 
    
      We finally note that for any scalars $a_1, \ldots , a_n$, we have 
      \[ \sum_{i = 1} ^n a_i \partial_i \vert _p  =  L_{\gamma} \] 
      with 
    \[  \gamma (t) : = \varphi ^{-1} (  ta_1, \ldots , ta_n) .   \]
    \end{proof}    
        
\end{document}